[\emph{Rambler}, July 1859.]

A Q\textsc{uestion} has arisen among
persons of theological knowledge and fair and candid minds, about the
wording and the sense of a passage in the \emph{Rambler} for May. It admits to my own mind of so clear and satisfactory
an explanation, that I should think it unnecessary to intrude myself, an
anonymous person, between the conductors and readers of this Magazine,
except that, as in dogmatic works the replies made to objections often
contain the richest matter, so here too, plain remarks on a plain
subject may open to the minds of others profitable thoughts, which are
more due to their own superior intelligence than to the very words of
the writer.

The \emph{Rambler},
then, has these words at p. 122: ``in the preparation of a
dogmatic definition, the faithful are consulted, as lately in the
instance of the Immaculate Conception.'' Now two questions bearing
upon doctrine have been raised on this sentence, putting aside the
question of fact as regards the particular instance cited, which must
follow the decision on the doctrinal questions: viz. first, whether it
can, with doctrinal correctness, be said that an \emph{appeal} to the
faithful is one of the preliminaries of a definition of doctrine; and
secondly, granting that the faithful are taken into account, still,
whether they can correctly be said to be \emph{consulted}. I shall
remark on both these points, and I shall begin with the second.

§ 1.

Now doubtless, if a divine were expressing himself
formally, and in Latin, he would not commonly speak of the laity being ``consulted''
among the preliminaries of a dogmatic definition, because the technical,
or even scientific, meaning of the word ``consult'' is to ``consult \emph{with},''
or to ``take \emph{counsel}.'' But the English word ``consult,'' in its
popular and ordinary use, is not so precise and narrow in its meaning;
it is doubtless a word expressive of trust and deference, but not of
submission. It includes the idea of inquiring into a matter of \emph{fact},
as well as asking a judgment. Thus we talk of ``consulting our barometer''
about the weather:—the barometer only attests the \emph{fact} of the
state of the atmosphere. In like manner, we may consult a watch or a
sun-dial about the time of day. A physician consults the pulse of his
patient; but not in the same sense in which his patient consults \emph{him}.
It is but an index of the state of his health. Ecclesiastes says, ``Qui \emph{observat}
ventum, non seminat;'' we might translate it, ``he who consults,'' without
meaning that we ask the wind's opinion. This being considered, it was, I
conceive, quite allowable for a writer, who was not teaching or treating
theology, but, as it were, conversing, to say, as in the passage in
question, ``In the preparation of a dogmatic definition, the faithful are
consulted.'' Doubtless their advice, their opinion, their judgment on the
question of definition is not asked; but the matter of fact, viz. their
belief, \emph{is} sought for, as a testimony to that apostolical
tradition, on which alone any doctrine whatsoever can be defined. In
like manner, we may ``consult'' the liturgies or the rites of the Church;
not that they speak, not that they can take any part whatever in the
definition, for they are documents or customs; but they are witnesses to
the antiquity or universality of the doctrines which they contain, and
about which they are ``consulted.'' And, in like manner, I certainly
understood the writer in the \emph{Rambler} to mean (and I think any lay
reader might so understand him) that the \emph{fidelium sensus} and \emph{consensus}
is a branch of evidence which it is natural or it necessary for
the Church to regard and consult, before she proceeds to any definition,
from its intrinsic cogency; and by consequence, that it ever has been so
regarded and consulted. And the writer's use of the word ``opinion'' in
the foregoing sentence, and his omission of it in the sentence in
question, seemed to show that, though the two cases put therein were
analogous, they were not identical.

Having said as much as this, I go further, and
maintain that the word ``consulted,'' used as it was used, was in no
respect unadvisable, except so far as it distressed any learned and good
men, who identified it with the Latin. I might, indeed, even have
defended the word as it was used, in the Latin sense of it. Regnier both
uses it of the laity and explains it. ``Cùm receptam apud populos
traditionem \emph{consulunt} et \emph{sequuntur} Episcopi, non illos
habent pro magistris et ducibus, \&c.'' (De Eccles. Christi. p. i. §
1, c. i., ed. Migne, col. 234.) But in my bountifulness I will give up
this use of the word as untheological; still I will maintain that the
true theological sense is unknown to all \emph{but} theologians.
Accordingly, the use of it in the \emph{Rambler} was in no sense
dangerous to any lay reader, who, if he knows Latin, still is not called
upon, in the structure of his religious ideas, to draw those careful
lines and those fine distinctions, which in theology itself are the very
means of anticipating and repelling heresy. The laity would not have a
truer, or a clearer, or a different view of the doctrine itself, though
the sentence had run, ``in the preparation of a dogmatic decree, \emph{regard}
is had to the sense of the faithful;'' or, ``there is an \emph{appeal} to
the general voice of the faithful;'' or, ``\emph{inquiry} is made into the
belief of the Christian people;'' or, ``the definition is not made without
a previous \emph{reference} to what the faithful will think of it and
say to it;'' or though any other form of words had been used, stronger or
weaker, expressive of the same general idea, viz. that \emph{the sense of
the faithful is not left out of the question} by the Holy See among
the preliminary acts of defining a doctrine.

Now I shall go on presently to remark on the
proposition itself which is conveyed in the words on which I have been
commenting; here, however, I will first observe, that such
misconceptions as I have been setting right will and must occur, from
the nature of the case, whenever we speak on theological subjects in the
vernacular; and if we do not use the vernacular, I do not see how the
bulk of the Catholic people are to be catechised or taught at all.
English has innovated on the Latin sense of its own Latin words; and if
we are to speak according to the conditions of the language, and
are to make ourselves intelligible to the multitude, we shall
necessarily run the risk of startling those who are resolved to act as
mere critics and scholastics in the process of popular instruction.

This divergence from a classical or ecclesiastical
standard is a great inconvenience, I grant; but we cannot remodel our
mother-tongue. \emph{Crimen} does not properly mean \emph{crime}; \emph{amiable}
does not yet convey the idea of \emph{amabilis}; \emph{compassio} is not
\emph{compassion}; \emph{princeps} is not a \emph{prince}; \emph{disputatio}
is not \emph{a dispute}; \emph{prævenire} is not \emph{to prevent}. \emph{Cicero
imperator} is not \emph{the Emperor Cicero}; \emph{scriptor egregius}
is not \emph{an egregious writer}; \emph{virgo singularis} is not \emph{a
singular virgin}; \emph{retractare dicta} is not \emph{to retract what
he has said}; and, as we know from the sacred passage, \emph{traducere}
is not necessarily \emph{to traduce}.

Now this is not merely sharp writing, for mistakes
do in matter of fact occur not unfrequently from this imperfect
correspondence between theological Latin and English; showing that
readers of English are bound ever to bear in mind that they are not
reading Latin, and that learned divines must ever exercise charity in
their interpretations of vernacular religious teaching.

For instance, I know of certain English sermons
which were translated into French by some French priests. They, good and
friendly men, were surprised to find in these compositions such language
as ``weak evidence and strong evidence,'' and ``insufficient, probable,
demonstrative evidence;'' they read that ``some writers had depreciated
the evidences of religion,'' and that ``the last century, when love was
cold, was an age of evidences.'' \emph{Evidentia}, they said, meant that
luminousness which attends on demonstration, conviction, certainty; how
can it be more or less? how can it be unsatisfactory? how can a sane man
disparage it? how can it be connected with religious coldness? The
simple explanation of the difficulty was, that the writer was writing
for his own people, and that in English ``an evidence'' is not \emph{evidentia}.

Another instance. An excellent Italian religious,
now gone to his reward, was reading a work of the same author; and he
came upon a sentence to the effect, I think, that the doctrine of the
Holy Trinity was to be held with ``implicit'' faith. He was perplexed and
concerned. He thought the writer held that the Church did not explicitly
teach, had not explicitly defined, the dogma; that is, he confused the
English meaning of the word, according to which it is a sort of
correlative to \emph{imperative}, meaning simple, unconditional,
absolute, with its sense in theology.

It is not so exactly apposite to refer,—yet I
will refer,— to another instance, as supplying a general
illustration of the point I am urging. It was in a third country that a
lecturer spoke in terms of disparagement of ``Natural Theology,'' on the
ground of its deciding questions of revelation by reasonings from
physical phenomena. It was objected to him, that \emph{Naturalis Theologia}
embraced \emph{all} truths and arguments from natural reason bearing
upon the Divine Being and Attributes. Certainly he would have been the
last to depreciate what he had ever made the paramount preliminary
science to Christian faith; but he spoke according to the sense of those
to whom his words might come. He considered that in the Protestant
school of Paley and other popular writers, the idea of Natural Theology
had practically merged in a scientific view of the argument from Design.

Once more. Supposing a person were to ask me
whether a friend, who has told me the fact in confidence, had written a
certain book, and I were to answer, ``Well, if he did, he certainly would
tell \emph{me},'' and the inquirer went away satisfied that he did \emph{not}
write it,—I do not see that I have done any thing to incur the
reproach of the English word ``equivocation;'' I have but adopted a mode
of turning-off a difficult question, to which any one may be obliged any
day to have recourse. I am not speaking of spontaneous and gratuitous
assertions, statements on solemn occasions, or answers to formal
authorities. I am speaking of impertinent or unjustifiable questions;
and I should like to know the man who thinks himself bound to say every
thing to every one. Physicians evade the questions of sick persons about
themselves; friends break bad news gradually, and with temporary
concealments, to those whom it may shock. Parents shuffle with their
children. Statesmen, ministers in Parliament, baffle adversaries in
every possible way short of a direct infringement of veracity. When St.
Athanasius saw that he was pursued on the Nile by the imperial officers,
he turned round his boat and met them; when they came up to his party
and hailed them, and asked whether they had seen any thing of
Athanasius, Athanasius cried out, ``O yes, he is not far from you;'' and
off the vessels went in different directions as swiftly as they could
go, each boat on its own errand, the pursuer and the pursued. I do not
see that there is in any of these instances what is expressed by the
English word ``equivocation;'' but it is the \emph{æquivocatio} of a
Latin treatise; and when Protestants hear that \emph{æquivocamus sine
scrupulo}, they are shocked at the notion of our ``unscrupulous
equivocation.''

Now, in saying all this, I must not be supposed to
be forgetful of the sacred and imperative duty of preserving with
religious exactness all those theological terms which are
ecclesiastically recognised as portions of dogmatic statements, such as \emph{Trinity},
\emph{Person}, \emph{Consubstantial}, \emph{Nature}, \emph{Transubstantiation},
\emph{Sacrament}, \&c. It would be unpardonable for a Catholic to
teach ``justification by faith only,'' and say that he meant by ``faith'' \emph{fides
formata}, or ``justification without works,'' and say that he meant by ``works''
the works of the Jewish ritual; but granting all this fully, still if
our whole religious phraseology is, as a matter of duty, to be modelled
in strict conformity to theological Latin, neither the poor nor children
will understand us. I have always fancied that to preachers great
license was allowed, not only in the wording, but even in the matter of
their discourses; they exaggerate and are rhetorical, and they are
understood \emph{piè} as speaking \emph{more prædicatorio}. I have
always fancied that, when Catholics were accused of hyperbolical
language towards the Blessed Virgin, it was replied that devotion was
not the measure of doctrine; nor surely is the vernacular of a magazine
writer. I do not see that I am wrong in considering that a periodical,
not treating theology \emph{ex professo}, but accidentally alluding to
an ecclesiastical act, commits no real offence if it uses an
unscientific word, since it speaks, not \emph{more digladiatorio}, but \emph{colloquialiter}.

I shall conclude this head of my subject with
allusion to a passage in the history of St. Dionysius the Great, Bishop
of Alexandria, though it is beyond my purpose; but I like to quote a
saint whom, \emph{multis nominibus} (not ``with many names,'' or `` by many
nouns''), I have always loved most of all the Ante-Nicene Fathers. It
relates to an attack which was made on his orthodoxy; a very serious
matter. Now I know every one will be particular on his own special
science or pursuits. I am the last man to find fault with such
particularity. Drill-sergeants think much of deportment; hard logicians
come down with a sledge-hammer even on a Plato who does not happen to
enumerate in his beautiful sentence all the argumentative considerations
which go to make up his conclusion; scholars are horrified, as if with
sensible pain, at the perpetration of a false quantity. I am far from
ridiculing, despising, or even undervaluing such precision; it is for
the good of every art and science that it should have vigilant
guardians. Nor am I comparing such precision (far from it) with that
true religious zeal which leads theologians to keep the sacred Ark of
the Covenant in every letter of its dogma, as a tremendous deposit for
which they are responsible. In this curious sceptical world, such
sensitiveness is the only human means by which the treasure of
faith can be kept inviolate. There is a woe in Scripture against the
unfaithful shepherd. We do not blame the watch-dog because he sometimes
flies at the wrong person. I conceive the force, the peremptoriness, the
sternness, with which the Holy See comes down upon the vagrant or the
robber, trespassing upon the enclosure of revealed truth, is the only
sufficient antagonist to the power and subtlety of the world, to
imperial comprehensiveness, monarchical selfishness, nationalism, the
liberalism of philosophy, the encroachments and usurpations of science.
I grant, I maintain all this; and after this avowal, lest I be
misunderstood, I venture to introduce my notice of St. Dionysius. He was
accused on a far worse charge, and before a far more formidable
tribunal, than commonly befalls a Catholic writer; for he was brought up
before the Holy See on a denial of our Lord's divinity. He had been
controverting with the Sabellians; and he was in consequence accused of
the doctrine to which Arius afterwards gave his name, that is, of
considering our Lord a creature. He says, writing in his defence, that
when he urged his opponents with the argument that ``a vine and a
vine-dresser were not the same,'' neither, therefore, were the ``Father
and the Son,'' these were not the only illustrations that he made use of,
nor those on which he dwelt, for he also spoke of ``a root and a plant,'' ``a
fount and a stream,'' which are not only \emph{distinct} from each other,
but of one and the same \emph{nature}. Then he adds, ``But my accusers
have no eyes to see this portion of my treatise; but they take up two
little words detached from the context, and proceed to discharge them at
me as pebbles from a sling.'' \footnote{Athan. de Sent. Dion. 8.}
If even a saint's words are not always precise enough to allow of being
made a dogmatic text, much less are those of any modern periodical.

The conclusion I would draw from all I have been
saying is this: Without deciding whether or not it is advisable to
introduce points of theology into popular works, and especially whether
it is advisable for laymen to do so, still, if this actually is done, we
are not to expect in them that perfect accuracy of expression which I
demanded in a Latin treatise or a lecture \emph{ex cathedrâ}; and if
there be a want of this exactness, we must not at once think it proceeds
from self-will and undutifulness in the writers.

§ 2.

Now I come to the \emph{matter} of what the writer
in the \emph{Rambler} really said, putting aside the question of the \emph{wording};
and I begin by expressing my belief that, whatever he may be
willing to admit on the score of theological Latinity in the use of the
word ``consult'' when applied to the faithful, yet one thing he cannot
deny, viz. that in using it, he implied, from the very force of the
term, that they are treated by the Holy See, on occasions such as that
specified, with attention and consideration.

Then follows the question, Why? and the answer is
plain, viz. because the body of the faithful is one of the witnesses to
the fact of the tradition of revealed doctrine, and because their \emph{consensus}
through Christendom is the voice of the Infallible Church.

I think I am right in saying that the tradition of
the Apostles, committed to the whole Church in its various constituents
and functions \emph{per modum unius}, manifests itself variously at
various times: sometimes by the mouth of the episcopacy, sometimes by
the doctors, sometimes by the people, sometimes by liturgies, rites,
ceremonies, and customs, by events, disputes, movements, and all those
other phenomena which are comprised under the name of history. It
follows that none of these channels of tradition may be treated with
disrespect; granting at the same time fully, that the gift of
discerning, discriminating, defining, promulgating, and enforcing any
portion of that tradition resides solely in the \emph{Ecclesia docens}.

One man will lay more stress on one aspect of
doctrine, another on another; for myself, I am accustomed to lay great
stress on the \emph{consensus fidelium}, and I will say how it has come
about.

1. It had long been to me a difficulty, that I
could not find certain portions of the defined doctrine of the Church in
ecclesiastical writers. I was at Rome in the year 1847; and then I had
the great advantage and honour of seeing Fathers Perrone and Passaglia,
and having various conversations with them on this point. The point of
difficulty was this, that up to the date of the definition of certain
articles of doctrine respectively, there was so very deficient evidence
from existing documents that Bishops, doctors, theologians, held them. I
do not mean to say that I expressed my difficulty in this formal shape;
but that what passed between us in such interviews as they were kind
enough to give me, ran into or impinged upon this question. Nor would I
ever dream of making them answerable for the impression which their
answers made on me; but, speaking simply on my own responsibility, I
should say that, while Father Passaglia seemed to maintain that the
Ante-Nicene writers were clear in their testimonies in behalf (\emph{e.g}.)
of the doctrines of the Holy Trinity and Justification, expressly
praising and making much of the Anglican Bishop Bull; rather Perrone, on
the other hand, not speaking, indeed, directly upon those particular
doctrines, but rather on such as I will presently introduce in his own
words, seemed to me to say ``\emph{transeat}'' to the alleged fact which
constituted the difficulty, and to lay a great stress on what he
considered t be the \emph{sensus} and \emph{consensus fidelium}, as a
compensation or whatever deficiency there might be of patristical
testimony in behalf of various points of the Catholic dogma.

2. I should have been led to fancy, perhaps, that
he was shaping his remarks in the direction in which he considered he
might be especially serviceable to myself, who had been accustomed to
account for the (supposed) phenomena in another way, had it not been for
his work on the Immaculate Conception, which I read the next year with
great interest, and which was passing through the press when I saw him.
I am glad to have this opportunity of expressing my gratitude and
attachment to a venerable man, who never grudged me his valuable time.

But now for his treatise, to which I have referred,
so far as it speaks of the \emph{sensus fidelium}, and of its bearing
upon the doctrine, of which his work treats and upon its definition.

(1.) He states the historical \emph{fact} of such \emph{sensus}.
Speaking of the ``Ecclesiæ sensus'' on the subject, he says that, though
the liturgies, of the Feast of the Conception ``satis apertè patefaciant
quid Ecelesia antiquitùs de hoc senserit argumento,'' yet it may be
worth while to add some direct remarks on the sense itself of the
Church. Then he says, ``Ex duplici fonte eum colligi posse arbitramur,
turn scilicet ex pastorum, \emph{tum ex fidelium} sese gerendi ratione''
(pp. 74, 75). Let it be observed, he not only joins together the \emph{pastores}
and \emph{fideles}, but contrasts them; I mean (for it will bear on what
is to follow), the ``faithful'' do not include the ``pastors.''

(2.) Next he goes on to describe the relation of
that \emph{sensus fidelium} to the \emph{sensus Ecclesiæ}. He says,
that to inquire into the sense of the Church on any question, is nothing
else but to investigate towards which side of it she has more inclined.
And the ``indicia et manifestationes hujus propensionis'' are her public
acts, liturgies, feasts, prayers, ``pastorum ac fidelium in unum veluti
conspiratio'' (p. 101). Again, at p. 109, joining together in one his
twofold consent of pastors and people, he speaks of the ``unanimis
pastorum ac \emph{fidelium} consensio ... per liturgias, per festa, per
euchologia, per fidei controversias, per conciones patefacta.''

(3.) These various ``indicia'' are also the \emph{instrumenta
traditionis}, and vary one with another in the evidence which they
give in favour of particular doctrines; so that the strength of one
makes up in a particular case for the deficiency of another, and the
strength of the ``sensus communis fidelium'' can make up (\emph{e.g}.) for
the silence of the Fathers. ``Istiusmodi instrumenta interdum simul
conjunctè conspirare possunt ad traditionem aliquam apostolicam atque
divinam patefaciendam, interdum vero seorsum ... Perperam nonnulli
solent ad inficiandam traditionis alicujus existentiam urgere silentium
Patrum ... quid enim si silentium istud alio pacto … compensetur?'' (p.
139). He instances this from St. Irenæus and Tertullian in the ``Successio
Episcoporum,'' who transmit the doctrines ``tum activi operâ ministerii,
tum usu et praxi, tum institutis ritibus ... adeò ut catholica atque
apostolica doctrina inoculata ... fuerit ... communi Ecclesiæ cœtui''
(p. 142).

(4.) He then goes on to speak directly of the force
of the ``sensus fidelium,'' as distinct (not separate) from the teaching
of their pastors. ``\emph{Præstantissimi} theologi maximam \emph{probandi}
vim huic communi sensui inesse \emph{uno ore} fatentur. Etenim Canus, 'In
quæstione fidei,' inquit, 'communis fidelis populi sensus haud levem
facit fidem''' (p. 143). He gives another passage from him in a note,
which he introduces with the words, ``Illud præclarè addit;'' what Canus
adds is, ``Quæro ex te, quando de rebus Christianæ fidei inter nos
contendimus, non de philosophiæ decretis, utrùm potius \emph{quærendum}
est, quid philosophi atque ethnici, an quid homines Christiani, et \emph{doctrinâ
et fide} instituti, sentiant?'' Now certainly ``quærere quid sentiant
homines doctrinâ et fide instituti,'' though not asking advice, is an
act implying not a little deference on the part of the persons
addressing towards the parties addressed.

Father Perrone continues, ``Gregorius verò de
Valentiâ fusius vim ejusmodi fidelium consensus evolvit. 'Est enim,'
inquit, 'in \emph{definitionibus fidei} habenda ratio, quoad fieri
potest, consensus fidelium.''' Here, again, ``habere rationem,'' to have
regard to, is an act of respect and consideration. However, Gregory
continues, ``Quoniam \emph{et ii} sanè, quatenus \emph{ex ipsis} constat
Ecclesia, sic \emph{Spiritu Sancto assistente}, divinas revelationes \emph{integrè
et purè conservant}, ut omnes illi quidem aberrare non possunt ...
Illud solùm contendo: Si quando de re aliquâ in materie \emph{religionis}
controversia [controversâ?] constaret fidelium omnium concordem esse
sententiam (solet autem id constare, vel ex ipsâ praxi alicujus cultûs
communiter apud christianos populos receptâ, vel ex \emph{scandalo et}
\emph{offensione communi}, quæ opinione aliquâ oritur, \&c.)
meritò \emph{posse et debere} Pontificem illâ niti, ut quæ esset \emph{Ecclesiæ
sententia infallibilis}'' (p. 144). Thus Gregory says that, in
controversy about a matter of faith, the consent of all the faithful has
such a force in the proof of this side or that, that the Supreme Pontiff
\emph{is able and ought} to \emph{rest} upon it, as being the \emph{judgment
or sentiment} of the \emph{infallible} Church. These are surely
exceedingly strong words; not that I take them to mean strictly that
infallibility is \emph{in} the ``consensus fidelium,'' but that that ``consensus''
is an \emph{indicium} or \emph{instrumentum} to us of the judgment of
that Church which \emph{is} infallible.

Father Perrone proceeds to quote from Petavius, who
supplies us with the following striking admonition from St. Paulinus,
viz. ``ut de omnium fidelium \emph{ore pendeamus}, quia in omnem fidelem
Spiritus Dei spirat.''

Petavius speaks thus, as he quotes him (p. 156): ``\emph{Movet
me}, ut in eam [viz. piam] sententiam sim propensior, \emph{communis
maximus sensus fidelium omnium}.'' By ``movet me'' he means, that he \emph{attends}
to what the \emph{cœtus fidelium} says: this is certainly not passing
over the fideles, but making much of them.

In a later part of his work (p. 186), Father
Perrone speaks of the ``consensus fidelium'' under the strong image of a \emph{seal}.
After mentioning various arguments in favour of the Immaculate
Conception, such as the testimony of so many universities, religious
bodies, theologians, \&c., he continues, ``Hæc demum omnia firmissimo
veluti \emph{sigillo obsignat} totius christiani populi consensus.''

(5.) He proceeds to give several instances, in
which the definition of doctrine was made in consequence of nothing else
but the ``sensus fidelium'' and the ``juge et vivum magisterium'' of the
Church.

For his meaning of the ``juge et vivum magisterium
Ecclesiæ,'' he refers us to his \emph{Prælectiones} (part ii. § 2, c.
ii.). In that passage I do not see that he defines the sense of the
word; but I understand him to mean that high authoritative voice or act
which is the Infallible Church's prerogative, inasmuch as she is the
teacher of the nations; and which is a sufficient warrant to all men for
a doctrine being true and being \emph{de fide}, by the mere fact of its
formally occurring. It is distinct from, and independent of, tradition,
though never in fact separated from it. He says, ``Fit ut traditio
dogmatica identificetur cum ipsâ Ecclesiæ doctninâ, a quâ separari
nequit; qua propter, \emph{etsi documenta deficerent omnia}, solum hoc
vivum et juge magisterium \emph{satis esset} ad cognoscendam doctrinam
divinitus traditam, habito præsertim respectu ad solennes Christi
promissiones'' (p. 303).

This being understood, he speaks of several points
of faith which have been determined and defined by the ``magisterium'' of
the Church and, as to tradition, on the ``consensus fidelium,''
prominently, if not solely.

The most remarkable of these is the ``dogma de
visione Dei beatificâ'' possessed by souls after purgatory and before
the day of judgment; a point which Protestants, availing themselves of
the comment of the Benedictines of St. Maur upon St. Ambrose, are
accustomed to urge in controversy. ``Nemo est qui nesciat,'' says Father
Perrone, ``quot utriusque Ecclesiæ, tum Græcæ tum Latinæ, Patres
contrarium sensisse visi sunt'' (p. 147). He quotes in a note the words
of the Benedictine editor, as follows: ``Propemodum incredibile videri
potest, quàm in eâ quæstione sancti Patres ab ipsis Apostolorum
temporibus ad Gregorii XI. [Benedicti XII.] pontificatum florentinumque
concilium, hoc est toto quatuordecim seculorum spatio, incerti ac parùm
constantes exstiterint.'' Father Perrone continues: ``Certè quidem in
Ecclesiâ non deerat quoad hunc fidei articulum divina traditio;
alioquin nunquam is definiri potuisset: verùm non omnibus illa erat
comperta; divina eloquia haud satis in re sunt. conspicua; \emph{Patres},
ut vidimus, in varias abierunt sententias; \emph{liturgiæ ipsæ} non
modicam præ se ferunt difficultatem. \emph{His omnibus sucurrit} juge
Ecclesiæ magisterium, \emph{communis præterea fidelium sensus}; qui
altè adeò defixum … habebant mentibus, purgatas animas statim ad
Deum videndum eoque fruendum admitti, ut non minimum eorum animi vel ex
ipsâ controversiâ fuerint \emph{offensi}, quæ sub Joanne XXII.
agitabatur, et cujus definitio \emph{diu nimis protrahebatur}.'' Now does
not this imply that the tradition, on which the definition was made, was
manifested in the \emph{consensus fidelium} with a luminousness which
the succession of Bishops, though many of them were ``Sancti Patres ab
ipsis Apostolorum temporibus,'' did not furnish? that the definition was
delayed till the \emph{fideles} would bear the delay no longer? that it
was made because of them and for their sake, because of their strong
feelings? If so, surely, in plain English, most considerable deference
was paid to the ``sensus fidelium;'' their opinion and advice indeed was
not asked, but their testimony was taken, their feelings consulted,
their impatience, I had almost said, feared.

In like manner, as regards the doctrine, though not
the definition, of the Immaculate Conception, he says, not denying, of
course, the availableness of the other ``instrumenta traditionis'' in this
particular case, ``Ratissimum est, Christi fideles omnes circa hunc
articulum unius esse animi, idque ita, ut maximo afficerentur \emph{scandalo},
si vel minima de Immaculatâ Virginis Conceptione quæstio
moveretur'' (p. 156).

3. A year had hardly passed from the appearance of
Fr. Perrone's book in England, when the Pope published his Encyclical
Letter. In it he asked the Bishops of the Catholic world, ``ut nobis
significare velitis, quâ devotione vester clerus \emph{populusque fidelis}
erga Immaculatæ Virginis conceptionem sit animatus, et quo desiderio
flagret, ut ejusmodi res ab apostolicâ sede decernatur;'' that is, when
it came to the point to take measures for the definition of the
doctrine, he did lay a special stress on this particular preliminary,
viz. the ascertainment of the feeling of the faithful both towards the
doctrine and its definition; as the \emph{Rambler} stated in the passage
out of which this argument has arisen. It seems to me important to keep
this in view, whatever becomes of the word ``consulted,'' which, I have
already said, is not to be taken in its ordinary Latin sense.

4. At length, in 1854, the definition took place,
and the Pope's Bull containing it made its appearance. In it the Holy
Father speaks as he had spoken in his Encyclical, viz. that although he \emph{already}
knew the sentiments of the Bishops, still he had wished to know the
sentiments of the \emph{people} also: ``\emph{Quamvis} nobis ex receptis
postulationibus de definiendâ tandem aliquando Immaculatâ Virginis
conceptione perspectus esset plurimorum sociorum \emph{Antistitum}
sensus, tamen Encyclicas literas, \&c. ad omnes Ven. FF. totius
Catholici orbis sacrorum Antistites misimus, ut, adhibitis ad Deum
precibus, nobis scripto \emph{etiam} significarent, quæ esset suorum \emph{fidelium}
erga Immaculatam Deiparæ Conceptionem pietas et devotio,'' \&c. And
when, before the formal definition, he enumerates the various witnesses
to the apostolicity of the doctrine, he sets down ``divina eloquia,
veneranda traditio, perpetuus Ecclesiæ sensus, singularis catholicorum
Antistitum ac \emph{fidelium} conspiratio.'' \emph{Conspiratio}, the two,
the Church teaching and the Church taught, are put together, as one
twofold testimony, illustrating each other, and never to be divided.

5. A year or two passed, and the Bishop of
Birmingham published his treatise on the doetrine. I close this portion
of my paper with an extract from his careful view of the argument. ``Nor
should the universal conviction of pious Catholics be passed over, as of
small account in the general argument; for that pious belief, and the
devotion which springs from it, are the \emph{faithful reflection} of
the pastoral teaching'' (p. 172). Reflection; that is, the people are a \emph{mirror},
in which the Bishops see themselves. Well, I suppose a person may \emph{consult}
his glass, and in that way may know things about himself which he
can learn in no other way. This is what Fr. Perrone above seems to say
has sometimes actually been the case, as in the instance of the ``beatifica
visio'' of the saints; at least he does not mention the ``\emph{pastorum}
ac fidelium \emph{conspiratio}'' in reviewing the grounds of its
definition, but simply the ``juge Ecclesiæ magisterium'' and the ``communis
fidelium sensus.''

His lordship proceeds: ``The more devout the
faithful grew, the more devoted they showed themselves towards this
mystery. And it is the devout who have the surest instinct in discerning
the mysteries of which the Holy Spirit breathes the grace through the
Church, and who, with as sure a tact, reject what is alien from her
teaching. The common accord of the faithful has weighed much as an
argument even with the most learned divines. St. Augustine says, that
amongst many things which most justly held him in the bosom of the
Catholic Church, was the 'accord of populations and of nations.' In
another work he says, 'It seems that I have believed nothing but the
confirmed opinion and the exceedingly wide-spread report of populations
and of nations.' Elsewhere he says: 'In matters whereupon the Scripture
has not spoken clearly, the custom of the people of God, or the
institutions of our predecessors, are to be held as law.' In the same
spirit St. Jerome argues, whilst defending the use of relics against
Vigilantius: 'So the people of all the Churches who have gone out to
meet holy relics, and have received them with so much joy, are to be
accounted foolish''' (pp. 172, 173).

And here I might come to an end; but, having got so
far, I am induced, before concluding, to suggest an historical instance
of the same great principle, which Father Perrone does not draw out.

§ 3.

First, I will set down the various ways in which
theologians put before us the bearing of the Consent of the faithful
upon the manifestation of the tradition of the Church. Its \emph{consensus}
is to be regarded: 1. as a testimony to the fact of the apostolical
dogma; 2. as a sort of instinct, or [\emph{phron\={e}ma}], deep in
the bosom of the mystical body of Christ; 3. as a direction of the Holy
Ghost; 4. as an answer to its prayer; 5. as a jealousy of error, which
it at once feels as a scandal.

1. The first of these I need not enlarge upon, as
it is illustrated in the foregoing passages from Father Perrone.

2. The second is explained in the well-known
passages of Möhler's \emph{Symbolique}; e.g. ``L'esprit de Dieu, qui
gouverne et vivifie l'Eglise, enfante dans l'homme, en s'unissant
à lui, \emph{un instinct}, un tact éminemment chrétien, qui le
conduit à toute vraie doctrine ... Ce sentiment commun, cette
conscience de l'Eglise est la tradition dans le sens subjectif du mot.
Qu'est-ce donc que la tradition considérée sous ce point de vue? C'est
le sens chrétien existant dans l'Eglise, et transmis par l'Eglise; sens,
toutefois, qu'on ne peut séparer des vérités qu'il contient, puisqu'il
est formé de ces vérités et par ces vérités.'' Ap. Perrone, p. 142.

3. Cardinal Fisher seems to speak of the third, as
he is quoted by Petavius, \emph{De Incarn}. xiv. 2; that is, he speaks
of a custom imperceptibly gaining a position, ``nullâ præceptorum vi,
sed consensu quodam tacito tam populi quàm cleri, quasi tacitis omnium
suffragiis recepta fuit, priusquàm ullo conciliorum decreto legimus eam
fuisse firmatam.'' And then he adds, ``This custom has its birth \emph{in
that people which is ruled by the Holy Ghost},'' \&c.

4. Petavius speaks of a fourth aspect of it. ``It is
well said by St. Augustine, that to the minds of individuals certain
things are revealed by God, not only by extraordinary means, as in
visions, \&c., but also in those usual ways, according to which what
is unknown to them is opened \emph{in answer to their prayer}. After
this manner it to be believed, that God has revealed to Christians the
sinless Conception of the Immaculate Virgin.'' \emph{De Incarn}. xiv. 2,
11.

5. The fifth is ``enlarged upon in Dr. Newman's
second \emph{Lecture on Anglican Difficulties}, from which I quote a few
lines: ``We know that it is the property of life to be impatient of any
foreign substance in the body to which it belongs. It will be sovereign
in its o domain, and it conflicts with what it cannot assimilate into
itself, and \emph{is irritated and disordered} till it has expelled it.
Such expulsion, then, is emphatically a test of uncongeniality, for it
shows that the substance ejected, not only is not one with the body that
rejects it, but cannot be made one with it; that its introduction is not
only useless or superfluous, adventitious, but that it is intolerable.''
Presently he continues: ``The religious life of a people is of a certain
quality and direction, and these are tested by the mode in which it
encounters the various opinions, customs, and institutions which are
submitted to it. Drive a stake into a river's bed, and you will at once
ascertain which way it is running, and at what speed; throw up even a
straw upon the air, and you will see which way the wind blows; submit
your heretical and Catholic principle to the action of the multitude,
and you will be able to pronounce at once whether it is imbued with
Catholic truth or with heretical falsehood.'' And then he proceeds
to exemplify this by a passage in the history of Arianism, the very
history which I intend now to take, as illustrative of the truth and
importance of the thesis on which I am insisting.

It is not a little remarkable, that, though,
historically speaking, the fourth century is the age of doctors,
illustrated, as it was, by the saints Athanasius, Hilary, the two
Gregories, Basil, Chrysostom, Ambrose, Jerome, and Augustine, and all of
these saints bishops also, except one, nevertheless in that very day the
divine tradition committed to the infallible Church was proclaimed and
maintained far more by the faithful than by the Episcopate.

Here, of course, I must explain:—in saying this,
then, undoubtedly I am not denying that the great body of the Bishops
were in their internal belief orthodox; nor that there were numbers of
clergy who stood by the laity, and acted as their centres and guides;
nor that the laity actually received their faith, in the first instance,
from the Bishops and clergy; nor that some portions of the laity were
ignorant, and other portions at length corrupted, by the Arian teachers,
who got possession of the sees and ordained an heretical clergy;—but I
mean still, that in that time of immense confusion the divine dogma of
our Lord's divinity was proclaimed, enforced, maintained, and (humanly
speaking) preserved, far more by the ``Ecclesia docta'' than by the ``Ecclesia
docens;'' that the body of the episcopate was unfaithful to its
commission, while the body of the laity was faithful to its baptism;
that at one time the Pope, at other times the patriarchal, metropolitan,
and other great sees, at other times general councils, said what they
should not have said, or did what obscured and compromised revealed
truth; while, on the other hand, it was the Christian people who, under
Providence, were the ecclesiastical strength of Athanasius, Hilary,
Eusebius of Vercellæ, and other great solitary confessors, who would
have failed without them.

I see, then, in the Arian history a palmary example
of a state of the Church, during which, in order to know the tradition
of the Apostles, we must have recourse to the faithful; for I fairly
own, that if I go to writers, since I must adjust the letter of Justin,
Clement, and Hippolytus with the Nicene Doctors, I get confused; and
what revives and re-instates me, as far as history goes, is the faith of
the people. For I argue that, unless they had been catechised, as St.
Hilary says, in the orthodox faith from the time of their baptism, they
never could have had that horror, which they show, of the heterodox
Arian doctrine. Their voice, then, is the voice of tradition; and
the instance comes to us with still greater emphasis, when we
consider—1. that it occurs in the very beginning of the history of the
``Ecclesia docens,'' for there can scarcely be said to be any history of
her teaching till the age of martyrs was over; 2. that the doctrine in
controversy was so momentous, being the very foundation of the Christian
system; 3. that the state of controversy and disorder lasted over the
long space of sixty years; and that
it involved serious persecutions, in life, limb, and property, to the
faithful whose loyal perseverance decided it.

It seems, then, as striking an instance as I could
take in fulfilment of Father Perrone's statement, that the voice of
tradition may in certain cases express itself, not by Councils, nor
Fathers, nor Bishops, but the ``communis fidelium sensus.''

I shall set down some authorities for the two
points successively, which I have to enforce, viz. that the Nicene dogma
was maintained during the greater part of the 4th century,

1. not by the unswerving firmness of the Holy See,
Councils, or Bishops, but

2. by the ``consensus fidelium.''

I. On the one hand, then, I say, that there was a
temporary suspense of the functions of the ``Ecclesia docens.'' The body
of Bishops failed in the confession of the faith. They spoke variously,
one against another; there was nothing, after Nicæa, of firm,
unvarying, consistent testimony, for nearly sixty years. There were
untrustworthy Councils, unfaithful Bishops; there was weakness, fear of
consequences, misguidance, delusion, hallucination, endless, hopeless,
extending itself into nearly every corner of the Catholic Church. The
comparatively few who remained faithfu1 were discredited and driven into
exile; the rest were either deceivers or were deceived.

1. A.D.
325. The great council of Nicæa, of 318 Bishops, chiefly from the
eastern provinces of Christendom, under the presidency of Hosius of
Cordova, as the Pope's Legate. It was convoked against Arianism, which
it once for all anathematized; and it inserted the formula of the ``Consubstantial''
into the Creed, with the view of establishing the fundamental dogma
which Arianism impugned. It is the first Œcumenical Council, and
recognised at the time its own authority as the voice of the infallible
Church. It is so received by the \emph{orbis terrarum} at this day. The
history of the Arian controversy, from its date, A.D.
325, to the date of the second Œcumenical Council, A.D. 381, is the history of the struggle through Christendom for
the universal acceptance or the repudiation of the formula of the ``Consubstantial.''

2. A.D.
334, 335. The synods of Cæsarea and Tyre against Athanasius, who was
therein accused and formally condemned of rebellion, sedition, and
ecclesiastical tyranny; of murder, sacrilege, and magic; deposed from
his see, forbidden to set foot in Alexandria for life, and banished to
Gaul. Constantine confirmed the sentence.

3. A.D.
341. Council of Rome of fifty Bishops, attended by the exiles from
Thrace, Syria, \&c., by Athanasius, \&c., in which Athanasius was
pronounced innocent.

4. A.D.
341. Great Council of the Dedication at Antioch, attended by ninety or a
hundred Bishops. The council ratified the proceedings of the councils of
Cæsarea and Tyre, and placed an Arian in the see of Athanasius. Then it
proceeded to pass a dogmatic decree in reversal of the formula of the ``Consubstantial.''
Four or five creeds, instead of the Nicene, were successively adopted by
the assembled fathers. The first was a creed which they ascribed to
Lucian, a martyr and saint of the preceding century, in whom the Arians
always gloried as their master. The second was fuller and stronger in
its language, and made more pretension to orthodoxy. The third was more
feeble again. These three creeds were circulated in the neighbourhood;
but, as they wished to send one to Rome, they directed a fourth to be
drawn up. This, too, apparently failed. So little was known at the time
of the real history of this synod and its creeds, that St. Hilary calls
it ``sanctorum synodus.''

5. A.D.
345. Council of the creed called Macrostich. This creed suppresses, as
did the third, the word ``substance.'' The eastern Bishops sent this to
the Bishops of the West, who rejected it.

6. A.D.
347. The great council of Sardica, attended by 380 Bishops. Before it
commenced, the division between its members broke out on the question
whether or not Athanasius should have a seat in it. In consequence,
seventy-six retired to Philippopolis, on the Thracian side of Mount Hæmus,
and there excommunicated the Pope and the Sardican fathers. These
seceders published a sixth confession of faith. The synod of Sardica,
including Bishops from Italy, Gaul, Africa, Egypt, Cyprus, and
Palestine, confirmed the act of the Roman council, and restored
Athanasius and the other exiles to their sees. The synod of
Philippopolis, on the contrary, sent letters to the civil magistrates of
those cities, forbidding them to admit the exiles into them. The
imperial power took part with the Sardican fathers, and Athanasius went
back to Alexandria.

7. A.D.
351. Before many years had run out, the great eastern party was up
again. Under pretence of putting down a kind of Sabellianism, they drew
up a new creed, into which they introduced certain inadvisable
expressions of some of the ante-Nicene writers, on the subject of our
Lord's divinity, and dropped the word ``substance.'' St. Hilary thought
this creed also Catholic; and other Catholic writers style its fathers ``holy
Bishops.''

8. There is considerable confusion of dates here.
Anyhow, there was a second Sirmian creed, in which the eastern
party first came to a division among themselves. St. Hilary at length
gives up these creeds as indefensible, and calls this one a ``blasphemy.''
It is the first creed which criticises the words ``substance,'' \&c.,
as unscriptural. Some years afterwards this ``blasphemia'' seems to have
been interpolated, and sent into the East in the name of Hosius. At a
later date, there was a third Sirmian Creed; and a second edition of it,
with alterations, was published at Nice in Thrace.

9. A.D.
353. The council of Arles. I cannot find how many Bishops attended it.
As the Pope sent several Bishops as legates, it must have been one of
great importance. The Bishop of Arles was an Arian, and managed to
seduce, or to force, a number of orthodox Bishops, including the Pope's
legate, Vincent, to subscribe the condemnation of Athanasius. Paulinus,
Bishop of Trêves, was nearly the only champion of the Nicene faith and
of Athanasius. He was accordingly banished into Phrygia, where he died.

10. A.D.
355. The council of Milan, of more than 300 Bishops of the West. Nearly
all of them subscribed the condemnation of Athanasius; whether they
generally subscribed the heretical creed, which was brought forward,
does not appear. The Pope's four legates remained firm, and St.
Dionysius of Milan, who died an exile in Asia Minor. An Arian was put
into his see. Saturninus, the Bishop of Arles, proceeded to hold a
council at Beziers; and its fathers banished St. Hilary to Phrygia.

11. A.D.
357. Hosius falls. Constantius used such violence towards the old man,
and confined him so straitly, that at last, broken by suffering, he was
brought, though hardly, to hold communion with Valens and Ursacius [the
Arian leader], though he would not subscribe against Athanasius.'' Athan.
\emph{Arian. Hist}. 45.

12. Liberius. A.D. 357. ``The tragedy was not ended in the lapse of Hosius, but
in the evil which befell Liberius, the Roman Pontiff, it became far more
dreadful and mournful, considering that he was Bishop of so great a
city, and of the whole Catholic Church, and that he had so bravely
resisted Constantine two years previously. There is nothing, whether in
the historians and holy fathers, or in his own letters, to prevent our
coming to the conclusion, that Liberius communicated with the Arians,
and confirmed the sentence passed against Athanasius; but he is not at
all on that account to be called a heretic.'' Baron. Ann. 357, 40-45.
Athanasius says: ``Liberius, after he had been in banishment two years,
gave way, and from fear of threatened death was induced to subscribe.'' \emph{Arian.
Hist}. § 41. St. Jerome says: ``Liberius, tædio victus exilii, in hæreticam
pravitatem subscribens, Romam quasi victor intravit.'' \emph{Chron}.

13. A.D.
359. The great councils of Seleucia and Ariminum, being one bi-partite
council, representing the East and West respectively. At Seleucia there
were 150 Bishops, of which only the twelve or thirteen from Egypt were
champion of the Nicene ``Consubstantial.'' At Ariminum there were as ma y
as 400 Bishops, who, worn out by the artifice of long delay on the
part of the Arians, abandoned the ``Consubstantial,'' and subscribed the
ambiguous formula which the heretics had substituted for it.

14. A.D.
361. The death of Constantius; the Catholic Bishops breathe again, and
begin at once to remedy the miseries of the Church, though troubles were
soon to break out anew.

15. A.D.
362. State of the Church of Antioch at this time. There were four
Bishops or communions of Antioch; first, the old succession and
communion, which had possession before the Arian troubles; secondly, the
Arian succession, which had lately conformed to orthodoxy in the person
of Meletius; thirdly, the new Latin succession, lately created by
Lucifer, whom some have thought the Pope's legate there; and, fourthly,
the new Arian succession, which was begun upon the recantation of
Meletius. At length, as Arianism was brought under, the evil reduced
itself to two successions, that of Meletius and the Latin, which went on
for many years, the West and Egypt holding communion with the latter,
and the East with the former.

16. A.D.
370-379. St. Basil was Bishop of Cæsarea in Cappadocia through these
years. The judgments formed about this great doctor in his lifetime show
us vividly the extreme confusion which prevailed. He was accused by one
party of being a follower of Apollinaris, and lost in consequence some
of the sees over which he was metropolitan. He was accused by the monks
in his friend Gregory's diocese of favouring the semi-Arians. He was
accused by the Neocæsareans of inclining towards Arianism. And he was
treated with suspicion and coldness by Pope Damasus.

17. About A.D. 360, St. Hilary says: ``I am not speaking of things foreign
to my knowledge; I am not writing about what I am ignorant of; I have
heard and I have seen the shortcomings of persons who are present to me,
not of laymen merely, but of Bishops. For, excepting the Bishop Eleusius
and a few with him, for the most part the ten Asian provinces, within
whose boundaries I am situate, are truly ignorant of God.'' It is
observable, that even Eleusius, who is here spoken of as somewhat better
than the rest, was a semi-Arian, according to Socrates, and even a
persecutor of Catholics at Constantinople; and, according to Sozomen,
one of those who urged Pope Liberius to give up the Nicene formula of
the ``Consubstantial.'' By the ten Asian provinces is meant the east and
south provinces of Asia Minor, pretty nearly as cut off by a line
passing from Cyzicus to Seleucia through Synnada.

18. A.D.
360. St. Gregory Nazianzen says, about this date: ``Surely the pastors
have done foolishly; for, excepting a very few, who, either on account
of their insignificance were passed over, or who by reason of their
virtue resisted, and who were to be left as a seed and root for the
springing up again and revival of Israel by the influences of the
Spirit, all temporised, only differing from each other in this, that
some succumbed earlier, and others later; some were foremost champions
and leaders in the impiety, and others joined the second rank of
the battle, being overcome by fear, or by interest, or by flattery, or,
what was the most excusable, by their own ignorance.'' \emph{Orat}. xxi.
24.

19. A.D.
363. About this time, St. Jerome says: ``Nearly all the churches in the
whole world, under the pretence of peace and the emperor, are polluted
with the communion of the Arians.'' \emph{Chron}. Of the same date, that
is, upon the Council of Ariminum, are his famous words, ``Ingemuit totus
orbis et se esse Arianum miratus est.'' \emph{In Lucif}. That is, the
Catholics of Christendom were surprised indeed to find that their rulers
had made Arians of them.

20. A.D.
364. And St. Hilary: ``Up to this date, the only cause why Christ's
people is not murdered by the priests of Anti-christ, with this deceit
of impiety, is, that they take the words, which the heretics use, to
denote the faith which they themselves hold. Sanctiores aures plebis quàm
corda sunt sacerdotum.'' \emph{In Aux}. 6.

21. St. Hilary speaks of the series of
ecclesiastical councils of that time in the following well-known
passage: ``It is most dangerous to us, and it is lamentable, that there
are at present as many creeds as there are sentiments, and as many
doctrines among us as dispositions, while we write creeds and explain
them according to our fancy. Since the Nicene council, we have done
nothing but write the creed. While we fight about words, inquire about
novelties, take advantage of ambiguities, criticise authors, fight on
party questions, have difficulties in agreeing, and prepare to
anathematise each other, there is scarce a man who belongs to Christ.
Take, for instance, last year's creed, what alteration is there not in
it already? First, we have the creed, which bids us not to use the
Nicene 'consubstantial;' then comes another, which decrees and preaches
it; next, the third, excuses the word 'substance,' as adopted by the
fathers in their simplicity; lastly, the fourth, instead of excusing,
condemns. We impose creeds by the year or by the month, we change our
minds about our own imposition of them, then we prohibit our changes,
then we anathematise our prohibitions. Thus, we either condemn others in
our own persons, or ourselves in the instance of others, and while we
bite and devour one another, are like to be consumed one of another.''

22. A.D.
382. St. Gregory writes: ``If I must speak the truth, I feel disposed to
shun every conference of Bishops; for never saw I synod brought to a
happy issue, and remedying, and not rather aggravating, existing evils.
For rivalry and ambition are stronger than reason,—do not think me
extravagant for saying so,—and a mediator is more likely to incur some
imputation himself than to clear up the imputations which others lie
under.'' \emph{Ep}. 129. It must ever be kept in mind that a passage like
this only relates, and is here quoted as only relating, to that
miserable time of which it is spoken. Nothing more can be argued from it
than that the ``Ecclesia docens'' is not at every time the active
instrument of the Church's infallibility.

II. Now we come secondly to the proofs of the fidelity of the laity, and
the effectiveness of that fidelity, during that domination of imperial
heresy to which the foregoing passages have related. I have abridged the
extracts which follow, but not, I hope, to the injury of their sense.

1. ALEXANDRIA.
``We suppose,'' says Athanasius, ``you are not ignorant what outrages they
[the Arian Bishops] committed at Alexandria, for they are reported every
where. They attacked \emph{the holy virgins and brethren} with naked
swords; they beat with scourges their persons, esteemed honourable in
God's sight, so that their feet were lamed by the stripes, whose souls
were whole and sound in purity and all good works.'' Athan. \emph{Op. c.
Arian}. 15, Oxf. tr.

``Accordingly Constantius writes letters, and
commences \emph{a persecution against all}. Gathering together a
multitude of herdsmen and shepherds, and dissolute youths belonging to
the town, armed with swords and clubs, they attacked in a body \emph{the
Church} of Quirinus: and \emph{some} they slew, \emph{some} they
trampled under foot, \emph{others} they beat with stripes and cast into
prison or banished. They haled away many \emph{women} also, and dragged
them openly into the court, and insulted them, dragging them by the
hair. \emph{Some} they proscribed; from \emph{some} they took away their
bread, for no other reason but that they might be induced to join the
Arians, and receive Gregory [the Arian Bishop], who had been sent by the
Emperor.'' Athan. \emph{Hist. Arian}. § 10.

``On the week that succeeded the holy Pentecost,
when the \emph{people}, after their fast, had gone out to the cemetery
to pray, because that \emph{all} refused communion with George [the
Arian Bishop], the commander, Sebastian, straightway with a multitude of
soldiers proceeded to \emph{attack the people}, though it was the Lord's
day; and finding a few praying, (for the greater part had already
retired on account of the lateness of the hour,) having lighted a pile,
he placed certain \emph{virgins} near the fire, and endeavoured to force
them to say that they were of the Arian faith. And having seized on \emph{forty
men}, he cut some fresh twigs of the palm-tree, with the thorns upon
them, and scourged them on the back so severely that some of them were
for a long time under medical treatment, on account of the thorns which
had entered their flesh, and others, unable to bear up under their
sufferings, died. All those whom they had taken, both the men and the
virgins, they sent away into banishment to the great oasis. Moreover,
they immediately banished out of Egypt and Libya the following Bishops
[sixteen], and the presbyters, Hierax and Dioscorus: some of them died
on the way, others in the place of their banishment. They caused also
more than thirty Bishops to take to flight.'' \emph{Apol. de Fug}. 7.

2. EGYPT.
``The Emperor Valens having issued an edict commanding that the orthodox
should be expelled both from Alexandria and the rest of Egypt, \emph{depopulation
and ruin to am immense extent immediately followed}; some were
dragged before the tribunals, others cast into prison, and many
tortured in various ways; all sorts of punishment being inflicted upon
persons who aimed only at peace and quiet.'' Socr. \emph{Hist}. iv. 24,
Bohn.

3. THE
MONKS OF
EGYPT. ``\emph{Antony
left the solitude} of the desert to go about every part of the city
[Alexandria], warning the inhabitants that the Arians were opposing the
truth, and that the doctrines of the Apostles were preached only by
Athanasius.'' Theod. \emph{Hist}. iv. 27, Bohn.

``Lucius, the Arian, with a considerable body of
troops, proceeded to the \emph{monasteries} of Egypt, where he in person
assailed the assemblage of holy men with greater fury than the ruthless
soldiery. When these excellent persons remained unmoved by all the
violence, in despair he advised the military chief to send the fathers
of the monks, the Egyptian Macarius and his namesake of Alexandria, into
exile.'' Socr. iv. 24.

OF
CONSTANTINOPLE.
``\emph{Isaac}, on seeing the emperor depart at the head of his army,
exclaimed, 'You who have declared war against God cannot gain His aid.
Cease from fighting against Him, and He will terminate the war. Restore
the pastors to their flocks, and then you will obtain a bloodless
victory.'' \emph{Ibid}. 34.

OF
SYRIA,
\&c. ``That these heretical doctrines [Apollinarian and Eunomian] did
not finally become predominant is \emph{mainly to be attributed to the zeal
of the monks} of this period; for \emph{all the monks} of Syria,
Cappadocia, and the neighboring provinces \emph{were sincerely attached to
the Nicene faith}. The same fate awaited them which had been
experienced by the Arians; for they incurred the full weight of the
popular odium and aversion, when it was observed that their sentiments
were regarded with suspicion by the monks.'' Sozom. \emph{Hist}. vii. 27,
Bohn.

OF
CAPPADOCIA.
``Gregory, the father of Gregory Theologus, otherwise a most excellent
man and a zealous defender of the true and Catholic religion, not being
on his guard against the artifices of the Arians, such was his
simplicity, received with kindness certain men who were contaminated
with the poison, and subscribed an impious proposition of theirs. This
moved the monks to such indignation, that they \emph{withdrew forthwith
from his communion}, and took with them, after their example, a \emph{considerable
part of his flock}.'' Ed. Bened. \emph{Monit. in Greg. Naz}. Orat. 6.

4. SYRIA.
``Syria and the neighbouring provinces were plunged into confusion and
disorder, for the Arians were very numerous in these parts, and had
possession of the churches. The members of the Catholic Church \emph{were
not}, however, \emph{few in numbers}. It was through their
instrumentality that the Church of Antioch was preserved from the
encroachments of the Arians, and enabled to resist the power of Valens.
Indeed, it appears that all the Churches which were governed by men who
were firmly attached to the faith did not deviate from the form of
doctrine which they had originally embraced.'' Sozom. vi. 21.

5. ANTIOCH.
``Whereas he (the Bishop Leontius) took part in the blasphemy of
Arius, he made a point of concealing this disease, partly \emph{for fear of
the multitude}, partly for the menaces of Constantius; so those who
followed the apostolical dogmas gained from him neither patronage nor
ordination, but these who held Arianism were allowed the fullest liberty
of speech, and were placed in the ranks of the sacred ministry. But
Flavian and Diodorus, who had embraced the ascetical life, and
maintained the apostolical dogmas, \emph{openly withstood} Leontius's
machinations against religious doctrine. They threatened that they would
retire from the communion of his Church, and would go to the West, and
reveal his intrigues. Though they were not as yet in the sacred
ministry, but were \emph{in the ranks of the laity}, night and day they
used to excite all the people to zeal for religion. They were the first
to divide the singers into two choirs, and to teach them to sing
alternately the strains of David. They too, assembling the devout at the
shrines of the martyrs, passed the whole night there in hymns to God.
These things Leontius seeing, did not think it safe to hinder them, for
he saw that \emph{the multitude was especially well affected} towards
those excellent persons. Nothing, however, could persuade Leontius to
correct his wickedness. It follows, that among the clergy were many who
were infected with the heresy: but \emph{the mass of the people were
champions of orthodoxy}.'' Theodor. \emph{Hist}. ii. 24.

6. EDESSA.
``There is in that city a magnificent church, dedicated to St. Thomas the
Apostle, wherein, on account of the sanctity of the place, religious
assemblies are continually held. The Emperor Valens wished to inspect
this edifice; when, having learned that \emph{all who usually congregated
there were opposed to the heresy} which he favoured, he is said to
have struck the prefect with his own hand, because he had neglected to
expel them thence. The prefect, to prevent the slaughter of \emph{so great
a number} of persons, privately warned them against resorting
thither. But his admonitions and menaces were alike unheeded; for on the
following day \emph{they all crowded to the church}. When the prefect
was going towards it with a large military force, a poor woman, leading
her own little child by the hand, hurried hastily by on her way to the
church, breaking through the ranks of the soldiery. The prefect,
irritated at this, ordered her to be brought to him, and thus addressed
her: 'Wretched woman, whither are you running in so disorderly a manner?'
She replied, 'To the same place that others are hastening.' 'Have you
not heard,' said he, 'that the prefect is about to put to death all that
shall be found there?' 'Yes,' said the woman, 'and therefore I hasten,
that I may be found there.' 'And whither are you dragging that little
child?' said the prefect. The woman answered, '\emph{That he also may be
vouchsafed the honour of martyrdom}.' The prefect went back and
informed the emperor that \emph{all were ready to die in behalf of their
own faith}; and added that it would be preposterous to destroy so
many persons at one time, and thus succeeded in restraining the emperor's
wrath.'' Socr. iv. 18. ``Thus was the Christian faith confessed by the \emph{whole
city} of Edessa.'' Sozom. vi. 18.

7. SAMOSATA.
``The Arians, having deprived this exemplary flock of their shepherd,
elected in his place an individual with whom \emph{none of the inhabitants
of the city}, whether poor or rich, servants or mechanics, husbandmen
or gardeners, men or women, young or old, would hold communion. \emph{He
was left quite alone}; no one even calling to see him, or exchanging
a word with him. It is, however, said that his disposition was extremely
gentle; and this is proved by what I am about to relate. One day, when
he went to bathe in the public baths, the attendants closed the doors;
but he ordered the doors to be thrown open, that the people might be
admitted to bathe with himself. Perceiving that they remained in a
standing posture before him, imagining that great deference towards
himself was the cause of this conduct, he arose and left the bath. \emph{These
people believed that the water had been contaminated by his heresy},
and ordered it to be let out and fresh water to be supplied. When he
heard of this circumstance, he left the city, thinking that he ought no
longer to remain in a place \emph{where he was the object of public
aversion and hatred}. Upon this retirement of Eunonius, Lucius was
elected as his successor by the Arians. Some young persons were amusing
themselves with playing at ball in the market-place; Lucius was passing
by at the time, and the ball happened to fall beneath the feet of the
ass on which he was mounted. \emph{The youths uttered loud exclamations,
believing that the bait was contaminated}. They lighted a fire, and
hurled the ball through it, believing that by this process the ball
would be purified. Although this was only a childish deed, and although
it exhibits the remains of ancient superstition, yet it \emph{is sufficient
to show the odium which the Arian faction had incurred in this city}.
Lucius was far from imitating the mildness of Eunonius, and he persuaded
the heads of government to exile most of the clergy.'' Theodor. iv. 15.

8. OSROENE.
``Arianism met with similar opposition at the same period in Osroëne and
Cappadocia. Basil Bishop of Cæsarea, and Gregory Bishop of Nazianzus,
were held in high admiration and esteem \emph{throughout these regions}.''
Sozom. vi. 21.

9. CAPPADOCIA.
``Valens, in passing through Cappadocia, did all in his power to injure
the orthodox, and to deliver up the churches to the Arians. He thought
to accomplish his designs more easily on account of a dispute which was
then pending between Basil and Eusebius, who governed the Church of Cæsarea.
This dissension had been the cause of Basil's departing to Pontus. \emph{The
people, and some of the most powerful and wisest men of the city},
began to regard Eusebius with suspicion, and to meditate a secession
from his communion. The emperor and the Arian Bishops regarded the
absence of Basil, and the hatred of the people towards Eusebius, as
circumstances that would tend greatly to the success of their designs. \emph{But
their expectations were utterly frustrated}. On the first
intelligence of the intention of the emperor to pass through Cappadocia,
Basil returned to Cæsarea, where he effected a reconciliation with
Eusebius. The projects of Valens were thus defeated, and he
returned with his Bishops.'' Sozom. vi. 19.

10. PONTUS.
``It is said that when Eulalius, Bishop of Amasia in Pontus, returned
from exile, he found that his Church had passed into the hands of an
Arian, and that \emph{scarcely fifty inhabitants of the city} had
submitted to the control of their new Bishop.'' Sozom. vii. 2.

11. ARMENIA.
``That company of Arians who came with Eustathius to Nicopolis had
promised that they would bring over this city to compliance with the
commands of the imperial vicar. This city had great ecclesiastical
importance, both because it was the metropolis of Armenia, and because
it had been ennobled by the blood of martyrs, and governed hitherto by
Bishops of great reputation, and thus, as Basil calls it, was the nurse
of religion and the metropolis of sound doctrine. Fronto, one of the
city presbyters, who had hitherto shown himself as a champion of the
truth, through ambition gave himself up to the enemies of Christ, and
purchased the bishopric of the Arians at the price of renouncing the
Catholic faith. This wicked proceeding of Eustathius and the Arians
brought a new glory instead of evil to the Nicopolitans, since it gave
them an opportunity of defending the faith. Fronto, indeed, the Arians
consecrated, \emph{but there was a remarkable unanimity of clergy and
people in rejecting him}. Scarcely one or two clerks sided with him;
on the contrary, he became the execration of all Armenia.'' \emph{Vita S.
Basil. Maurin}. pp. clvii. clviii.

12. NICOMEDIA.
``Eighty pious clergy proceeded to Nicomedia, and there presented to the
emperor a supplicatory petition complaining of the ill-usage to which
they had been subjected. Valens, dissembling his displeasure in their
presence, gave Modestus, the prefect, a secret order to apprehend these
persons and put them to death. The prefect, \emph{fearing that he should
excite the populace to a seditious movement} against himself, if he
attempted the public execution of so many, pretended to send them away
into exile,'' \&c. Socr. iv. 16.

13. ASIA
MINOR.
St. Basil says, about the year 372: ``Religious people keep silence, but
every blaspheming tongue is let loose. Sacred things are profaned; \emph{those
of the laity} who are sound in faith \emph{avoid the places of worship}
as schools of impiety, and raise their hands in solitude, with groans
and tears, to the Lord in heaven.'' \emph{Ep}. 93. Four years after he
writes: ``Matters have come to this pass; \emph{the people have left their
houses of prayer}, and assemble in deserts: a pitiable sight; \emph{women
and children, old men, and others infirm}, wretchedly faring in the
open air, amid the most profuse rains and snow-storms, and winds, and
frost of winter; and again in summer under a scorching sun. To this they
submit, because they \emph{will have no part in the wicked Arian leaven}.''
\emph{Ep}. 342. Again: ``Only one offence is now vigorously punished, an
accurate observance of our fathers' traditions. For this cause the pious
are driven from their countries, and transported into deserts. The \emph{people are in lamentations}, in continual tears at home and
abroad. There is a cry in the city, a cry in the country, in the roads,
in the deserts. Joy and spiritual cheerfulness are no more; our feasts
are turned into mourning; our houses of prayer are shut up, our altars
deprived of the spiritual worship.'' \emph{Ep}. 343.

14. SCYTHIA.
``There are in this country a great number of cities, of towns, and of
fortresses. According to an ancient custom which still prevails, all the
churches of the whole country are under the sway of one Bishop. Valens
[the emperor] repaired to the church, and strove to gain over the Bishop
to the heresy of Arius; but this latter manfully opposed his arguments,
and, after a courageous defence of the Nicene doctrines, quitted the
emperor, and proceeded to another church, \emph{whither he was followed by
the people. Valens was extremely offended at being left alone} in a
church with his attendants, and, in resentment, condemned Vetranio [the
Bishop] to banishment. Not long after, however, he recalled him,
because, I believe, \emph{he apprehended an insurrection}.'' Sozom. vi.
21.

15. CONSTANTINOPLE.
``Those who acknowledged the doctrine of consubstantiality were not only
expelled from the churches, but also from the cities. But although
expulsion at first satisfied them [the Arians], they soon proceeded to
the worse extremity of inducing compulsory communion with them, caring
little for such a desecration of the churches. They resorted to all
kinds of scourgings, a variety of tortures, and confiscation of
property. Many were punished with exile, some died under the torture,
and others were put to death while being driven from their country. \emph{These
atrocities were exercised throughout all the eastern cities}, but
especially at Constantinople.'' Socr. ii. 27.

The following passage is quoted for the substantial
fact which it contains, viz. the testimony of popular tradition to the
Catholic doctrine: ``At this period a union was nearly effected between
the Novatian and Catholic Churches; for, as they both \emph{held the same
sentiments concerning the Divinity}, and were subjected to a common
persecution, the members of both Churches assembled and prayed together.
The Catholics then possessed no houses of prayer, for the Arians had
wrested them from them.'' Sozom. iv. 20.

16. ILLYRIA.
``The parents of Theodosius were Christians, and were attached to the
Nicene doctrine, hence he took pleasure in the ministration of Ascholius
[Bishop of Thessalonica]. He also rejoiced at finding that the \emph{Arian
heresy had not been received in Illyria}.'' Sozom. vii. 4.

17. NEIGHBOURHOOD
OF MACEDONIA.
``Theodosius inquired concerning the religious sentiments which were
prevalent in the other provinces, and ascertained that, as far as
Macedonia, \emph{one form of belief was universally predominant},''
\&c. Ibid.

18. ROME.
``With respect to doctrine no dissension arose either at Rome or in any
other of the Western Churches. \emph{The people unanimously adhered to the
form of belief established at Nicæa}.'' Sozom. vi. 23.

``Not long after, Liberius (the Pope) was recalled
and re-instated in his see; for the people of Rome, \emph{having raised a
sedition, and expelled Felix} [whom the Arian party had intruded]
from their Church, Constantius deemed it \emph{inexpedient to provoke the
popular fury}.'' Socr. ii. 37.

``Liberius, returning to Rome, found the mind of the
mass of men alienated front him, because he had so shamefully yielded to
Constantius. And thus it came to pass, that those persons who had
hitherto kept aloof from Felix [the rival Pope], and had avoided his
communion in favour of Liberius, on hearing what had happened, \emph{left
him for Felix}, who raised the Catholic standard. Among others,
Damasus [afterwards Pope] took the side of Felix. Such had been, even
from the times of the Apostles, \emph{the love of Catholic discipline in
the Roman people}.'' Baron. ann. 357. He tells us besides, that the
people would not even go to the public baths, lest they should bathe
with the party of Liberius.

19. MILAN.
``At the council of Milan, Eusebius of Vercellæ, when it was proposed to
draw up a declaration against Athanasius, said that the council ought
first to be sure of the faith of the Bishops attending it, for he had
found out that some of them were polluted with heresy. Accordingly he
brought before the Fathers the Nicene creed, and said he was willing to
comply with all their demands, after they had subscribed that
confession. Dionysius, Bishop of Milan, at once took up the paper and
began to write his assent; but Valens [the Arian] violently pulled pen
and paper out of his hands, crying out that such a course of proceeding
was impossible. Whereupon, after much tumult, \emph{the question came
before the people, and great was the distress of all of them}; the
faith of the Church was impugned by the Bishops. \emph{They then, dreading
the judgment of the people}, transfer their meeting from the church
to the imperial palace.'' Hilar. \emph{in Const}. i.

``As the feast of Easter approached, the empress
sent to St. Ambrose to ask a church of him, where the Arians who
attended her might meet together. He replied, that a Bishop could not
give up the temple of God. The pretorian prefect came into the church,
where St. Ambrose was, \emph{attended by the people}, and endeavoured to
persuade him to yield up at least the Portian Basilica. \emph{The people
were clamorous against the proposal}; and the prefect retired to
report how matters stood to the emperor. The Sunday following, St.
Ambrose was explaining the creed, when he was informed that the officers
were hanging up the imperial hangings in the Portian Basilica, and that
upon this news the people were repairing thither. While he was offering
up the holy sacrifice, a second message came that the \emph{people had
seized an Arian priest} as he was passing through the street. He
despatched a number of his clergy to the spot to \emph{rescue the Arian
from his danger}. The court looked on this resistance of the people
as seditious, and immediately laid considerable fines upon \emph{the whole
body of the tradesmen} of the city. Several were thrown into prison.
In three days' time these tradesmen were fined two hundred pounds
weight of gold, and they said \emph{that they were ready to give as much
again, on condition that they might retain their faith}. The prisons
were filled with tradesmen: \emph{all the officers of the household},
secretaries, agents of the emperor, and dependent officers who served
under various counts, were kept within doors, and were forbidden to
appear in public under pretence that they should bear no part in the
sedition. \emph{Men of higher rank were menaced with severe consequences},
unless the Basilica were surrendered …

``Next morning the Basilica was surrounded by
soldiers; but it was reported, that \emph{these soldiers had sent to the
emperor to tell him} that if he wished to come abroad he might, and
that they would attend him, if he was going to the assembly of the
Catholics; otherwise, that they \emph{would go to that which would be held
by St. Ambrose}. Indeed, the solders were all Catholics, as well as
the citizens of Milan; there were no heretics there, except a few
officers of the emperor and. some Goths …

``St. Ambrose was continuing his discourse when he
was told that the emperor had withdrawn the soldiers from the Basilica,
and that he had restored to the tradesmen the fines which he had exacted
from them. \emph{This news gave joy to the people}, who expressed their
delight with applauses and thanksgivings; \emph{the soldiers themselves
were eager to bring the news}, throwing themselves on the altars, and
kissing them in token of peace.'' Fleury's \emph{Hist}. xviii. 41, 42,
Oxf. trans.

20. THE
SOLDIERY.
Soldiers having been mentioned in the foregoing extract, I add the
following passage. ``Terentius, a general distinguished by his valour and
by his piety, was able, on his return from Armenia, to erect trophies of
victory. Valens promised to give him every thing that he might desire.
But he asked not for gold or silver, for lands, power, or honours; \emph{he
requested that a church might be given to those who preached the
apostolical doctrines}.'' Theodor. iv. 32.

``Valens sent Trajan, the general, against the
barbarians. Trajan was defeated, and, on his return, the emperor
reproached him severely, and accused him of weakness and cowardice. But
Trajan replied with great boldness, ``It is not I, O emperor, who have
been defeated; for you, \emph{by fighting against God, have thrown the
barbarians upon His protection}. Do you not know who those are whom
you have driven from the churches, and who are those to whom you have
given them up? Arintheus and Victor, the other commanders, \emph{accorded
in what he had said}, and brought the emperor to reflect on the truth
of their remonstrances.'' Ibid. 33.

21. CHRISTENDOM
GENERALLY.
St. Hilary to Constantius: ``Not only in words, but in tears, we beseech
you to save the Catholic Churches from any longer continuance of these
most grievous injuries, and of their present intolerable persecutions
and insults, which moreover they are enduring, which is monstrous, from
our brethren. Surely your clemency should listen to the \emph{voice of
those who cry out so loudly}, 'I am a Catholic, I have no wish to be
a heretic.' It should seem equitable to your sanctity, most
glorious Augustus, that they who fear the Lord God and His judgment
should not be polluted and contaminated with execrable blasphemies, but \emph{should
have liberty to follow those Bishops and prelates} who observe
inviolate the laws of charity, and who desire a perpetual and sincere
peace. It is impossible, it is unreasonable, to mix true and false, to
confuse light and darkness, and bring into a union, of whatever kind,
night and day. \emph{Give permission to the populations to hear the
teaching of the pastors whom they have wished}, whom they fixed on,
whom they have chosen, to attend their celebration of the divine
mysteries, to offer prayers through them for your safety and prosperity.''
\emph{In Const}. i.

Now I know quite well what will be said to so elaborate a collection of
instances as I have been making. The ``lector benevolus'' will quote
against me the words of Cicero; ``Utitur in re non dubiâ testibus non
necessariis.'' This is sure to befall a man, when he directs the
attention of a friend to any truth which hitherto he has thought little
of. At first, he seems to be hazarding a paradox, and at length to be
committing a truism. The hearer is first of all startled, and then
disappointed; he ends by asking, ``Is this all?'' It is a curious
phenomenon in the philosophy of the human mind, that we often do not
know whether we hold a point or not, though we hold it; but when our
attention is once drawn to it, then forthwith we find it so much part of
ourselves, that we cannot recollect when we began to hold it, and we
conclude (with truth), and we declare, that it has always been our
belief. Now it strikes me as worth noticing, that, though Father Perrone
is so clear upon the point of doctrine which I have been urging in 1847,
yet in 1842, which is the date of my own copy of his \emph{Prælectiones},
he has not given the \emph{consensus fidelium} any distinct place in his
\emph{Loci Theologici}, though he has even given ``heretici'' a place
there. Among the \emph{Media Traditionis}, he enumerates the magisterium
of the Church, the Acts of the Martyrs, the Liturgy, usages and rites of
worship, the Fathers, heretics, Church history; but not a word, that I
can find, directly and separately, about the \emph{sensus fidelium}.
This is the more remarkable, because, speaking of the \emph{Acta Martyrum},
he gives a reason for the force of the testimony of the martyrs which
belongs quite as fully to the faithful generally; viz. that, as not
being theologians, they can only repeat that objective truth, which, on
the other hand, Fathers and theologians do but present subjectively, and
thereby coloured with their own mental peculiarities. ``We learn from
them,'' he says, ``what was the traditionary doctrine in both domestic and
public assemblies of the Church, without any admixture of private
and (so to say) subjective explanation, such as at times creates a
difficulty in ascertaining the real meaning of the Fathers; and so much
the more, because many of them were either women or ordinary and
untaught laymen, who brought out and avowed just what they believed in a
straightforward inartificial way.'' May we not conjecture that the
argument from the Consent of the Faithful was but dimly written among
the \emph{Loci} on the tablets of his intellect, till the necessities,
or rather the requirements, of the contemplated definition of the
Immaculate Conception brought the argument before him with great force?
Yet who will therefore for an instant suppose that he did not always
hold it? Perhaps I have overlooked some passage of his treatises, and am
in consequence interpreting his course of thought wrongly; but, at any
rate, what I seem to see in him, is what actually does occur from time
to time in myself and others. A man holds an opinion or a truth, yet
without holding it with a simple consciousness and a direct recognition;
and thus, though he has never denied, he has never gone so far as to
profess it.

As to the particular doctrine to which I have here
been directing my view, and the passage in history by which I have been
illustrating it, I am not supposing that such times as the Arian will
ever come again. As to the present, certainly, if there ever was an age
which might dispense with the testimony of the faithful, and leave the
maintenance of the truth to the pastors of the Church, it is the age in
which we live. Never was the Episcopate of Christendom so devoted to the
Holy See, so religious, so earnest in the discharge of its special
duties, so little disposed to innovate, so superior to the temptation of
theological sophistry. And perhaps this is the reason why the ``consensus
fidelium'' has, in the minds of many, fallen into the background. Yet
each constituent portion of the Church has its proper functions, and no
portion can safely be neglected. Though the laity be but the reflection
or echo of the clergy in matters of faith, yet there is something in the
``pastorum et fidelium \emph{conspiratio},'' which is not in the pastors
alone. The history of the definition of the Immaculate Conception shows
us this; and it will be one among the blessings which the Holy Mother,
who is the subject of it, will gain for us, in repayment of the
definition, that by that very definition we are all reminded of the part
which the laity have had in the preliminaries of its promulgation. Pope
Pius has given us a pattern, in his manner of defining, of the duty of
considering the sentiments of the laity upon a point of tradition, in
spite of whatever fullness of evidence the Bishops had already
thrown upon it.

In most cases when a definition is contemplated,
the laity will have a testimony to give; but if ever there be an
instance when they ought to be consulted, it is in the case of doctrines
which bear directly upon devotional sentiments. Such is the Immaculate
Conception, of which the \emph{Rambler} was speaking in the sentence
which has occasioned these remarks. The faithful people have ever a
special function in regard to those doctrinal truths which relate to the
Objects of worship. Hence it is, that, while the Councils of the fourth
century were traitors to our Lord's divinity, the laity vehemently
protested against its impugners. Hence it is, that, in a later age, when
the learned Benedictines of Germany and France were perplexed in their
enunciation of the doctrine of the Real Presence, Paschasius was
supported by the faithful in his maintenance of it. The saints, again,
are the object of a religious \emph{cultus}; and therefore it was the
faithful, again, who urged on the Holy See, in the time of John XXII.,
to declare their beatitude in heaven, though so many Fathers spoke
variously. And the Blessed Virgin is preeminently an object of devotion;
and therefore it is, I repeat, that though Bishops had already spoken in
favour of her absolute sinlessness, the Pope was not content without
knowing the feelings of the faithful.

Father Dalgairns gives us another case in point;
and with his words I conclude: ``While devotion in the shape of a dogma
issues from the high places of the Church, in the shape of devotion ...
it starts from below ... Place yourselves, in imagination, in a vast
city of the East in the fifth century. Ephesus, the capital of Asia
Minor, is all in commotion; for a council is to be held there, and
Bishops are flocking in from all parts of the world. There is anxiety
painted on every face; so that you may easily see that the question is
one of general interest ... Ask the very children in the streets what is
the matter; they will tell you that wicked men are coming to make out
that their own mother is not the Mother of God. And so, during a
live-long day of June, they crowd around the gates of the old
cathedral-church of St. Mary, and watch with anxious faces each Bishop
as he goes in. Well might they be anxious; for it is well known that
Nestorius has won the court over to his side. It was only the other day
that he entered the town, with banners displayed and trumpets sounding,
surrounded by the glittering files of the emperor's body-guard, with
Count Candidianus, their general and his own partisan, at their
head. Besides which, it is known for certain, that at least eighty-four
Bishops are ready to vote with him; and who knows how many more? He is
himself the patriarch of Constantinople, the rival of Rome, the imperial
city of the East; and then John of Antioch is hourly expected with his
quota of votes; and he, the patriarch of the see next in influence to
that of Nestorius, is, if not a heretic, at least of that wretched party
which, in ecclesiastical disputes, ever hovers between the two camps of
the devil and of God. The day wears on, and still nothing issues from
the church; it proves, at least, that there is a difference of opinion;
and as the shades of evening close around them, the weary watchers grow
more anxious still. At length the great gates of the Basilica are thrown
open; and oh, w hat a cry of joy bursts from the assembled crowd, as it
is announced to them that Mary has been proclaimed to be, what every one
with a Catholic heart knew that she was before, the Mother of God! …
Men, women, and children, the noble and the low-born, the stately matron
and the modest maiden, all crowd round the Bishops with acclamations.
They will not leave them; they accompany them to their homes with long
procession of lighted torches; they burn incense before them, after the
eastern fashion, to do them honour. There was but little sleep in
Ephesus that night; for very joy they remained awake: the whole town was
one blaze of light, for each window was illuminated.'' \footnote{Sacred Heart.}

My own drift is somewhat different from that which
has dictated this glowing description; but the substance of the argument
of each of us is one and the same. I think certainly that the \emph{Ecclesia
docens} is more happy when she has such enthusiastic partisans about
her as are here represented, than when she cuts off the faithful from
the study of her divine doctrines and the sympathy of her divine
contemplations, and requires from them \emph{fides implicita} in her
word, which in the educated classes will terminate in indifference, and
in the poorer in superstition.

O.
